\title{MegaByte: Predicting Million-byte Sequences with Multiscale Transformers}

\begin{abstract}
Autoregressive transformers are spectacular models for short sequences but scale poorly to long sequences such as high-resolution images, podcasts, code, or books. We propose \textsc{MegaByte}, a multi-scale decoder architecture that enables end-to-end differentiable modeling of sequences of over one million bytes.
\textsc{MegaByte} segments sequences into patches and uses a \emph{local} submodel within patches and a \emph{global} model between patches. This enables sub-quadratic self-attention, much larger feedforward layers for the same compute, and improved parallelism during decoding---unlocking  better performance at reduced cost for both training and generation. Extensive experiments show that \textsc{MegaByte} allows byte-level models to perform competitively with subword models on long context language modeling, achieve state-of-the-art density estimation on ImageNet, and model audio from raw files. Together, these results establish the  viability of tokenization-free autoregressive sequence modeling at scale.
\end{abstract}

\section{Introduction}
\label{sec:intro}
Sequences of millions of bytes are ubiquitous; for example, music, image, or video files typically consist of multiple megabytes. However, large transformer decoders (LLMs) typically only use several thousand tokens of context \citep{gpt3,Zhang2022OPT:Models}---both because of the quadratic cost of self-attention but also, more importantly, the cost of large feedforward networks per-position. This severely limits the set of tasks where LLMs can be applied.

We introduce \textsc{MegaByte}, a new approach to modeling long byte sequences. First, byte sequences are segmented into fixed-sized patches, loosely analogous to tokens. Our model then consists of three parts: (1) a \emph{patch embedder}, which simply encodes a patch by losslessly concatenating embeddings of each byte, (2) a \emph{global} module, a large autoregressive transformer that inputs and outputs patch representations and (3) a \emph{local} module, a small autoregressive model that predicts bytes within a patch. Crucially, we observe that for many tasks, most byte predictions are relatively easy (for example, completing a word given the first few characters), meaning that large networks per-byte are unnecessary, and a much smaller model can be used for intra-patch modelling.

The \textsc{MegaByte} architecture gives three major improvements over Transformers for long sequence modelling:

\begin{enumerate}
    \item \textbf{Sub-quadratic self-attention} Most work on long sequence models has focused on mitigating the quadratic cost of self-attention. \textsc{MegaByte} decomposes long sequences into two shorter sequences, and optimal patch sizes reduces the self-attention cost to $O(N^{\frac{4}{3}})$, which remains tractable for even long sequences.
    \item \textbf{Per-patch feedforward layers} In GPT3-size models, more than 98\% of FLOPS are used in computing position-wise feedforward layers. \textsc{MegaByte} uses large feedforward layers per-patch rather than per-position, enabling much larger and more expressive models for the same cost. With patch size $P$, where a baseline transformer would use the same feedforward layer with $m$ parameters $P$ times, \textsc{MegaByte} can use a layer with $mP$ parameters once for the same cost.
    \item \textbf{Parallelism in Decoding} Transformers must perform all computations serially during generation because the input to each timestep is the output from the previous timestep. By generating representations for patches in parallel, \textsc{MegaByte} allows greater parallelism during generation. For example, a \textsc{MegaByte} model with 1.5B parameters can generate sequences 40\% \emph{faster} than a standard 350M Transformer, whilst also improving perplexity when trained with the same compute. 
\end{enumerate}

Together, these improvements allow us to train much larger and better-performing models for the same compute budget, scale to very long sequences, and improve generation speed during deployment.

\textsc{MegaByte} also provides a strong contrast to existing autoregressive models that typically use some form of tokenization, where sequences of bytes are mapped to larger discrete tokens \citep{bpe,dalle,hubert}.

Tokenization complicates pre-processing, multi-modal modelling, and transfer to new domains, while hiding useful structure from the model. It also means that most state-of-the-art models are not truly end to end. The most widely used approaches to tokenization require language-specific heuristics \citep{gpt2} or lose information \citep{dalle}. Replacing tokenization with efficient and performant byte models would therefore have many advantages.

We conduct extensive experiments for both \textsc{MegaByte} and strong baselines. We use a fixed compute and data budget across all models to focus our comparisons solely on the model architecture rather than training resources, which are known to benefit all models. We find that \textsc{MegaByte} allows byte-level models to perform competitively with subword models on long context language modeling, achieve state-of-the-art perplexities for density estimation on ImageNet, and allow audio modelling from raw audio files. Together, these results establish the viability of tokenization-free autoregressive sequence modeling at scale.

\section{\textsc{MegaByte} Transformer}
\subsection{Overview}
\label{sec:fastslow}

\textsc{MegaByte} is an autoregressive model for efficiently modeling long input sequences. \textsc{MegaByte} is comprised of 3 components: (1) a \emph{patch embedder} that inputs a discrete sequence, embeds each element, and chunks it into patches of length $P$ (2) a large \emph{global} Transformer that contextualizes patch representations by performing self-attention over previous patches, and (3) a smaller \emph{local} Transformer that inputs a contextualized patch representation from the global model, and autoregressively predict the \emph{next} patch.

\subsection{Components}
$\textbf{Patch Embedder}$ with patch size of $P$ maps a byte sequence $x_{0..T}$ to a sequence of patch embeddings of length $K=\frac{T}{P}$ and dimension $P\cdot D_G$. 

First, each byte is embedded with a lookup table $E^{\text{global-embed}}\in\mathbb{R}^{V \times D_G}$ to an embedding of size $D_G$ and positional embeddings are added.

\begin{align}
h^\text{embed}_{t}&= E^\text{global-embed}_{x_t}+E^\text{pos}_t & t\in [0..T]
\end{align}

Then, byte embeddings are reshaped into a sequence of $K$ patch embeddings with dimension $P\cdot D_G$.

To allow autoregressive modelling, the patch sequence is padded to start with a trainable patch-sized padding embedding ($E^{\text{global-pad}} \in\mathbb{R}^{P\times D_G} $), and the last patch is removed from the input. This sequence is the input to the global model, and is denoted $h^\text{global-in}\in\mathbb{R}^{K\times (P\cdot D_G)}$.

\begin{align}
h^\text{global-in}_{k} = 
\begin{cases}
E^{\text{global-pad}} , & \text{if } k = 0, \\
h^\text{embed}_{((k-1)\cdot P):(k \cdot P)}, & k\in [1,..,K),
\end{cases}
\end{align}

$\textbf{Global Model}$ is a decoder-only Transformer with dimension $P\cdot D_G$ that operates on a sequence of $K$ patches. It incorporates a self-attention mechanism and causal masking to capture dependencies between patches. It inputs a sequence of $K$ patch representations $h^\text{global-in}_{0:K}$, and outputs an updated representation $h^\text{global-out}_{0:K}$ by performing self-attention over previous patches.

\begin{align}
h^\text{global-out}_{0:K} &= \text{transformer}^{\text{global}}(h^\text{global-in}_{0:K}) 
\end{align}

The output of the final global layer $h^\text{global}_{0:K}$ contains $K$ patch representations of dimension $P\cdot D_G$. For each of these, we reshape them into sequences of length $P$ and dimension $D_G$, where position $p$ uses dimensions $p \cdot D_G$ to $(p+1) \cdot D_G$. Each position is then projected to the dimension of the local model with a matrix $w^{\text{GL}}\in\mathbb{R}^{D_G\times D_L}$ where $D_L$ is the local model dimension. We then combine these with byte embeddings of size $D_L$ for the tokens in the \emph{next} patch $E^\text{local-embed}_{x_{(k\cdot P + p - 1)}}$. The local byte embeddings is offset by one with a trainable local padding embedding ($E^{\text{local-pad}} \in\mathbb{R}^{D_L}$) to allow autoregressive modelling within a patch. This results in a tensor $h^\text{local-in}\in\mathbb{R}^{K\times P \times D_L}$.

\begin{align}

h^\text{local-in}_{k,p} &= w^{\text{GL}}h^\text{global-out}_{k,(p\cdot D_G):((p+1)\cdot D_G)} + E^\text{local-embed}_{x_{(k\cdot P + p - 1)}}
\end{align}

$\textbf{Local Model}$  is a smaller decoder-only Transformer of dimension $D_L$ that operates on a single patch $k$ containing $P$ elements, each of which is the sum of an output from the global model and an embedding of the previous byte in the sequence. $K$ copies of the local models are run on each patch independently (and in parallel during training), computing a representation $h^\text{local-out}\in\mathbb{R}^{K\times P\cdot D_L} $.

\begin{align}
h^\text{local-out}_{k,0:P} &= \text{transformer}^{\text{local}}(h^\text{local-in}_{k,0:P}) 
\end{align}

Finally, we can compute the probability distribution over the vocabulary at each position. The $p$th element of the $k$th patch corresponds to element $t$ of the complete sequence, where $t=k\cdot P+p$:

\begin{align}
p(x_{t}|x_{0:t}) &= \text{softmax}{(E^\text{local-embed}h^\text{local-out}_{k,p})}_{x_t} 
\end{align}

\subsection{Variations and Extensions} We experiment with several extensions of \textsc{MegaByte}.

\subsubsection{Convolutional Patch Encoder}

One limitation of chunking sequences into patches is that it is not translation invariant, and byte sequences may receive a different representation depending on their position in the patch. This may mean, for example, that a model has to relearn the meaning of a word at different offsets. To mitigate this issue, we experimented with augmenting the Patch Embedder with causal convolutional layers, which allow translation-invariant contextual representations of the bytes before they are chunked into patches. We use a stack of convolutional layers, with filter sizes of 3, 5 and 7. 

\subsubsection{Cross-patch Attention} The Local model uses short sequences for efficiency, and relies on the Global model for long-range information. However, we can increase the context of the Local model with little overhead by allowing it to condition on $r$ elements from the previous patch. This approach allows the Global model to focus on a longer-range context. Specifically, when computing self-attention in each layer, we concatenate the keys and values with the last $r$ keys and queries from the previous patch. We use rotary embeddings \citep{rotary} to model relative positions between elements in the sequence. This approach is reminiscent of TransformerXL \citep{xfmxl} but differs by being fully differentiable. 

\subsubsection{Strided Inference}
\label{strided}
We observed empirically that the per-token loss within each patch would increase towards the end of the patch, as the prediction relies more on the weaker Local model. To alleviate this issue, we propose \emph{strided inference}, in which we predict the sequence with two forward passes of the full model, whose inputs are offset by $p/2$ positions from each other. We then combine the first $p/2$ positions in each patch for our predictions to predict the complete sequence. Similarly to sliding window techniques \citep{shortformer}, this approach doubles the cost of inference but improves results.

\subsection{Motivation}
Having described the model, we briefly discuss the motivation behind some of the architectural choices.

\textbf{Why is the local model needed?} Many of the efficiency advantages of the \textsc{MegaByte} design could be realized with the Global model alone, which would resemble a decoder version of ViT \citep{vit}. However,  the joint distribution over the patch $p(x_{t+1},..,x_{t+P}|x_{0..t})$ has an output space of size $256^{P}$ so direct modeling is only tractable for very small patches. We could instead factor the joint distribution into conditionally independent distributions $p(x_{t+1}|x_{0..t})..p(x_{t+P}|x_{0..t})$, but this would greatly limit the model's expressive power. For example, it would be unable to express a patch distribution such as 50\% \emph{cat} and 50\% \emph{dog}, and would instead have to assign probability mass to strings such as \emph{cag} and \emph{dot}. Instead, our autoregressive Local model conditions on previous characters within the patch, allowing it to only assign probability to the desired strings.

\textbf{Increasing Parameters for Fixed Compute}  Transformer models have shown consistent improvements with parameter counts \citep{kaplan2020scaling}. However, the size of models is limited by their increasing computational cost. \textsc{MegaByte} allows larger models for the same cost, both by making self attention sub-quadratic, and by using large feedforward layers across patches rather than individual tokens.

\textbf{Re-use of Established Components} \textsc{MegaByte} consists of two transformer models interleaved with shifting, reshaping and a linear projection. This re-use increases the likelihood that the architecture will inherit the desirable scaling properties of transformers.

\section{Efficiency Analysis}
\subsection{Training Efficiency} We analyze the cost of different architectures when scaling both the sequence length and size of the models.

\paragraph{Attention} The cost of the self attention in a transformer architecture for a sequence of length $T$ has $O(T^2)$ complexity. Much work has been explored reducing this; for example, Sparse Transformers \citep{sparsetransformer} and Routing Transformers \citep{routing} show strong results with a complexity $O(T^\frac{3}{2})$. Numerous linear attention mechanisms have also been proposed \citep{lineartransformer,linear, performer}, although we are not aware of competitive results on large scale language modeling tasks. As a function of sequence length $T$ and patch size $P$, the Global model has a sequence of length $\frac{P}{T}$ so uses $O(\frac{T^2}{P^2})$ operations, and the Local model uses $\frac{P}{T}$ sequences of length $P$ so uses $O(\frac{TP^2}{P})=O(PT)$ operations. The overall cost of \textsc{MegaByte} is therefore in $O(\frac{T^2}{P^2}+TP)$. $P$ is a hyperparameter that is chosen to create an architecture for sequences of size $T$. By setting $P=T^\frac{1}{3}$ the complexity is in $O(T^\frac{4}{3})$. Using much shorter patches of $P=T^\frac{1}{5}$ would give a complexity of $O(T^\frac{8}{5})$. The cost is less than the transformer for all non-trivial values of $P$ such that $1<P<T$. 

\paragraph{Feedforward Layers} However, attention is not the main cost in large transformers. Instead of increasing the sequence length, transformers are more commonly scaled by increasing the dimension of their latent state $d$, and the feedforward network cost dominates the model's overall cost \citep{kaplan2020scaling}. For example, in the GPT3 architecture, the quadratic self-attention computation accounts for only 1.4\% of FLOPS. Following the approximation of \citep{kaplan2020scaling}, a forward pass with a large transformer with $m$ non-embedding parameters on a sequence of length $T$ uses roughly $2mT$ FLOPS. \textsc{MegaByte} contains two transformers: the Global model uses $m_g$ parameters on a sequence of length $\frac{T}{P}$, and a Local model with $m_l$ parameters that sees $\frac{T}{P}$ sequences of length $P$, giving an estimate of $2T(\frac{m_g}{P} + m_l)$ FLOPS. When $m_g\gg m_l$, the FLOPS used by \textsc{MegaByte} is approximately $\frac{2Tm_g}{P}$, allowing a model $P$ times larger than a transformer with equivalent FLOPS. This analysis holds irrespective of any efficient attention mechanisms used in the transformer.

\paragraph{Combined Analysis} To understand efficiency at different sequence lengths and model sizes, we calculate the total FLOPS used by transformers, Linear Transformers and \textsc{MegaByte}. For each operation, we use FLOP estimates from \citep{kaplan2020scaling}, except for attention in Linear Transformers, which we estimate as $9D$ FLOPS/token\footnote{This may underestimate the time taken by Linear Transformer decoders, which use a recurrence mechanism that is harder to parallelize on current hardware.}, where $D$ is the model embedding dimension. Figure \ref{figure:flops} shows that for models of size 660M to 173B and sequence lengths of up to 1M tokens, \textsc{MegaByte} with $P=8$ uses less FLOPS than either transformers or Linear Transformers. Baseline model architectures are based on GPT3, and Megabyte global/local model sizes are 452M/151M, 5.8B/604M, 170B/3.2B respectively.

\subsection{Generation Efficiency} Generating long sequences with transformers is slow, because the input to each timestep is the output from the previous timestep, meaning each layer must be computed for each token serially. As running a layer on a single token typically does not saturate the amount of parallelism available within a GPU, for analysis, we model each layer as a constant cost independently of size. Consider a \textsc{MegaByte} model with $L_{\text{global}}$ layers in the Global model and $L_{\text{local}}$ layers in the Local model and patch size $P$, compared with a Transformer architecture with $L_{\text{local}}+L_{\text{global}}$ layers. Generating each patch with \textsc{MegaByte} requires a sequence of $O(L_{\text{global}}+P\cdot L_{\text{local}})$ serial operations, whereas the Transformer requires $O(P\cdot L_{\text{global}}+P\cdot L_{\text{local}})$ serial operations. When $L_{\text{global}}\gg L_{\text{local}}$ (i.e. the Global model has many more layers than the Local model), \textsc{MegaByte} can reduce inference costs by a factor close to $P$.

\section{Experimental setup}
\label{sec:exp}

\subsection{Controlling for Compute and Data}
    \label{sec:controlled_exp}
Models show consistent improvements when increasing both data and compute \cite{kaplan2020scaling,Chinchilla}, meaning that one model can outperform another because of an increased training budget instead of an improved architecture. However, in practice, both compute and data are typically limited. We conduct experiments using a fixed compute and data budget across all models to focus comparisons solely on the model architecture rather than training resources. To achieve this, we adjust model hyperparameters (mainly, number of layers) within each architecture so that the forward pass time taken per byte is matched, and then train all models for the same number of bytes.

\subsection{Comparison Systems} We compare \textsc{MegaByte} with both a standard decoder-only Transformer and PerceiverAR \citep{perceiverar}. PerceiverAR extends the original transformer with a single cross-attention layer over a much longer context sequence, and is the best performing general purpose autoregressive model we are aware of and achieves state-of-the-art results across several modalities. We implemented both models in the same codebase, and all models share a similar data loader, preprocessing step, and trainer to avoid any artifacts in our compute-controlled experiments.

\subsection{Training Procedure}
    \label{sec:training_proc}
    All models were trained using the Metaseq\footnote{https://github.com/facebookresearch/metaseq} code base~\cite{OPT}. The training used the PyTorch framework \cite{pytorch}, with fairscale to improve memory efficiency through fully sharded model and optimizer states \cite{fairscale}. Mixed precision training was used to improve training efficiency at scale \cite{mixedprecision}. More training details and various model parameters can be found in Section \ref{subsec:append_training} in the Appendix. 

To validate our implementation of PerceiverAR, we reproduced their experiments on downsized ImageNet at 64 pixels. By carefully matching hyperparameters, we achieved a bits per byte (bpb) score of 3.53, compared to the reported 3.54 in the original paper.

\subsection{Inference Methods} Several techniques have been proposed for trading off speed for performance during inference with language models, including sliding windows \cite{shortformer} and our strided inference (Section \ref{strided}). We only use these methods when comparing with prior published work (Tables \ref{tab:text_sota} and \ref{tab:img_sota}).

\section{Language Modeling} We evaluated the performance of \textsc{MegaByte} on language modeling on a set of 5 diverse datasets emphasizing long-range dependencies: Project Gutenberg (PG-19), Books, Stories, arXiv, and Code. 

\textbf{Datasets} We experiment on a range of long form text datasets. The PG-19 dataset \cite{pg19} consists of English-language books written before 1919 and is extracted from the \href{https://www.gutenberg.org/}{Project Gutenberg online library}. The Stories dataset \cite{stories} is a subset of CommonCrawl data meant to emulate Winograd schemas. Books \cite{pile} is another collection of English-language books. The arXiv dataset is a collection of technical publications written in \LaTeX ~from the \href{arxiv.org}{arXiv online archive}. Finally, the Code dataset is a large publicly available dataset of open source code, under Apache, BSD or MIT licenses. More details on dataset sizes and document lengths are shared in Table \ref{tab:document_lengths}. 

\textbf{Controlled Experiments} Table \ref{tab:text_exps}, lists bpb on each dataset. Each model is trained for 80 billion bytes, and models are scaled to use the same compute budget. We carefully tune hyperparameters for all architectures to best utilize the available compute budget. \textsc{MegaByte} consistently outperforms both baseline transformers and PerceiverAR across all datasets. We use the same set of parameters on all datasest. In all experiments presented in Table~\ref{tab:text_exps}, transformer has size of 320M with context length of 1024, PerceiverAR has size of 248M with context size of 8192 and latent size of 1024, and \textsc{MegaByte} global/local model sizes are 758M/262M with context length of 8192 and patch size of 8. 

\label{sec:results}

\textbf{Scaling Experiment} We scale up our training data on PG-19 (Table~\ref{tab:pg19_summary}), and compare \textsc{MegaByte} with byte baselines, as well as converting all results to word-level perplexities to benchmark with state-of-art token based models. 

We train a byte-level Transformer, PerceiverAR and \textsc{MegaByte} models for 400B bytes and the same compute budget using same model parameters as in the controlled experiments. We find that \textsc{MegaByte} outperforms other byte-level models by a wide margin at this scale.\footnote{The only prior byte-level experiments we are aware of are at a smaller scale in \citet{blockrecurrent}, who report results equivalent to test perplexities of 46.5 with a version of the BlockRecurrent transformer, and 49.5 with Memorizing Transformers \cite{memorizing}, compared to 36.4 with our model. } 

We also compare with the best previously reported numbers for sub-word models. These results may be confounded by differing amounts of compute and tuning used, but show that \textsc{MegaByte} 
gives results competitive with state-of-the-art models trained on subwords. These results suggest that \textsc{MegaByte} may allow future large language models to be tokenization-free.

\section{Image Modeling}
\label{sec:image}

\subsection{Sequence Modeling on ImageNet}

We test \textsc{MegaByte} on variants of the autoregressive image generation task on ImageNet \citep{Oord2016PixelNetworks}, to measure its ability to efficiently use long context. We test on three different resolutions of images, ranging from 64\texttimes64 to 640\texttimes640 pixels -- the latter requiring the effective modeling of sequences with over 1.2M tokens. This generation task becomes increasingly challenging as the image's resolution grows: doing well on this task requires the modeling of local patterns (textures, lines, etc.) and long-range context that provides information about the high level structure of the image. Inspired by recent works in Vision Transformers~\citep{vit}, we model image data patch by patch (more details can be found in Appendix~\ref{sec:scan_type}). 

\subsection{Comparison with State of the Art}

We train a large \textsc{MegaByte} model on ImageNet 64x64 with Global and Local models sized 2.7B and 350M parameters, respectively, for 1.4T tokens. We estimate that training this model consumed less than half the GPU hours we would have needed to reproduce the best PerceiverAR model described by \citep{perceiverar}. As shown in Table \ref{tab:img_sota}, \textsc{MegaByte} matches the state-of-the-art performance of PerceiverAR whilst using only half the compute.

\subsection{Scaling to higher resolutions}

We compare three transformer variants (vanilla, PerceiverAR, \textsc{MegaByte}) to test scalability to long sequences on increasingly large image resolutions. We use our own implementations of these in the same framework and budget the same amount of GPU hours and data to train each of these model variants. 

\textsc{MegaByte} is able to handle all sequence lengths with a single forward pass of up to 1.2M tokens. We found neither the standard Transformer nor PerceiverAR could model such long sequences at a reasonable model size, so instead we split images into segments of size 1024 and 12000 respectively. For Megabyte, we set patch size as 12 for Image64 and patch size as 192 for Image256 and Image640 datasets. Model sizes are adjusted to match overall training speeds across models and we do not use any form of sliding window evaluation in this experiment. As seen in Table \ref{tab:resolution_scale}, \textsc{MegaByte} outperforms baselines across all resolutions in this compute-controlled setting. The precise settings used for each of the baseline models  such as context length and number of latents are summarized in Table \ref{table:opt_model_desc}.

Results show that \textsc{MegaByte} outperforms the other systems at all resolutions, demonstrating an effective model of sequences of over 1M bytes.

\section{Audio Modeling}
\label{sec:speech}

Audio has aspects of both the sequential structure of text and the continuous nature of images, so is an interesting application for \textsc{MegaByte}.

Raw audio is typically stored as a sequence of 16-bit integer values (one per timestep); a softmax layer would need to output 65,536 probabilities per timestep to model all possible values. To address this issue, various techniques have been developed to reduce the memory and computational requirements of the softmax layer. For instance, \citet{wavenet} apply $\mu$-law companding transformation and quantizes the input into 256 possible values. Alternatively, \citet{parawavenet} model the samples using the discretized mixture of logistics distribution introduced by \citet{pixcelcnn++}. Finally, \citet{wavernn} use a dual softmax technique to produce 8 coarse and 8 fine bits. In our approach, we simplify the audio modeling process by directly reading the bytes (256 possible values) from the audio file and conducting an autoregressive language model on top of that. This greatly streamlines the modeling process, making it easier and more efficient.

Our audio modeling approach focuses on 16 kHz, 16-bit audio, which equates to 32k bytes per one-second clip. We use an extensive audio dataset consisting of 2 terabytes (roughly 18,000 hours) of audio. We use a sequence length of 524,288, a patch size of 32, and a batch size of 32 to facilitate model training. By utilizing these settings, we can effectively train our model on large volumes of audio data, helping to improve its accuracy and efficacy.

Our model obtains bpb of 3.477, much lower than the results with perceiverAR (3.543) and vanilla transformer model (3.567). More ablation results are presented in Table \ref{tab:ablations}. 

\section{Analysis}
\subsection{Generation speed}

We also compare the text generation speed between \textsc{MegaByte} and a transformer. We compare a 350M parameter baseline transfomer and a \textsc{MegaByte} model with a 1.3B parameter Global model and a 218M parameter local model, trained on PG19 with equal compute. As shown in Table~\ref{tab:generation}, the \textsc{MegaByte} model achieves much lower perplexity as expected. However, \textsc{MegaByte} also generates a sequence of 8192 tokens 40\% \emph{faster} than transformer, despite having over 4 times the parameters. This speed up is due to the bulk of the parameters being in the Global model, which only needs to be computed once for every 8 tokens, whereas all the parameters in the baseline model are used on every token.

\subsection{Model Components} In Table \ref{tab:ablations}, we  analyze the significance of different components in the \textsc{MegaByte} architecture by studying arXiv, Librilight-L and ImageNet256 datasets. 

Removing Local (\textit{w/o local model}) or global (\textit{w/o global model}) model, we observe a substantial increase in bpb on all datasets, showing that both parts are crucial. The performance of the model without the cross-patch local model (\textit{w/o cross-patch local model}) is competitive, indicating that the architecture is robust to this modification. We observe slight improvement on the Librilight-L and ImageNet256 datasets by augmenting the \textsc{MegaByte} model with a CNN encoder (\textit{w/ CNN encoder}). This suggests that the \textsc{MegaByte} architecture can benefit from integrating alternative encoding mechanisms.

\subsection{Effective Use of Context} Long-context models often struggle to benefit from the full context \citep{sun2021long}. Figure \ref{fig:prob_within_context} shows that later tokens within each context window consistently have a higher likelihood, indicating that \textsc{MegaByte} can effectively use at least 8k bytes of context on the PG19 dataset. 

\subsection{Strided Inference}
\label{strided_inference}
 We find that within a single patch, on average, the \textsc{MegaByte} performs worse on later tokens within a patch (see Figure \ref{fig:strided_inference}). 

 Section \ref{strided} proposes \emph{strided inference} as a solution, where two forward passes are performed offset by $\frac{P}{2}$ tokens, and results from the first half of each patch are combined. Table \ref{tab:inference_options} shows performance improvements from strided inference, which are additive with the standard sliding window.

\subsection{Hyperparameters}

\textsc{MegaByte} introduces several additional hyperparameters. We tuned these parameters independently for different modalities and reported performance based on the best setting we found. All experiments in the same group use the same compute. 

\textbf{Patch Size.} We experimented with various patch sizes on Image256 dataset and found that there is a wide range of values where \textsc{MegaByte} performs similarly. We found similar robustness against the choice of this hyperparameter across all modalities, although the optimal patch size itself can be different across modalities.

\textbf{Local to Global model Size Ratio.} We experimented with different Local/Global model size ratios on PG19 dataset. By grouping bytes into patches, \textsc{MegaByte} effectively uses $P$ times less tokens for the Global model as on the Local model---enabling us to increase the size of the Global model without reduced cost. We find that a given compute budget is spent optimally when the Global model has more parameters than the Local model. This trend was consistent across all modalities and various patch sizes.

\section{Related Work}
\label{sec:related}
Prior research has explored the possibility of improving the efficiency of Transformers on long sequences, primarily motivated by mitigating the quadratic cost of self-attention.

\textbf{Efficient Encoder Models} Several related techniques to ours have been developed for transformer encoder architectures but cannot be straightforwardly applied to decoders. In particular, patchifying operations have previously been used in image \emph{encoder} models such as ViT \citep{vit}, and down- and up-sampling operations have been used for text encoders \citep{canine}, but such methods cannot be naively applied to decoder-only models without leaking information to future bytes in the same patch. \textsc{MegaByte} generalizes these approaches to an efficient decoder model by using a intra-patch transformer to predict each sequence element's likelihood, and offseting the inputs to the two models to avoid leaking information. \citet{perceiver} which uses self-attention on a shorter latent sequence, and \citet{temporalbottleneck} which uses recurrent model to process chunks with $k$ input steps also resemble patchification, but this technique cannot easily be applied to decoder architectures without leaking information to future timesteps.

\textbf{Efficient Decoder models} Improving the efficiency of decoder models is more challenging because of the need to make one prediction per timestep, and not leak information to future timesteps. The most popular approaches can be categorized as (1) chunking sequences into smaller blocks, and propagating information from previous blocks with either recurrence \citep{xfmxl,blockrecurrent} or cross-attention \citep{perceiverar}, (2) linear alternatives to attention, which typically involve forms of token-level recurrence \citep{lineartransformer, linear} or state space models \citep{s4,s5,mega}, or (3) sparse approximations of attention \citep{reformer,longformer,sparsetransformer,memorizing}. However, the performance of dense attention means it is typically still chosen for large scale decoders \citep{llama,palm}. \textsc{MegaByte} takes the alternative approach of decomposing the complete sequence into two shorter sequences, giving sub-quadratic attention. We also note that feedforward networks are the dominant cost in large decoders, not self-attention. Our approach to compressing sequences allows much larger models than would be possible when using large feedforward networks at every timestep.

\textbf{Tokenization} The most common approach to shortening sequence lengths in Transformer decoders is to pre-process the input with a form of tokenization, in which multiple bytes are mapped to a single discrete token from a fixed vocabulary. For text, this can be done losslessly using methods such as BPE \citep{bpe} and SentencePiece \citep{sentencepiece}, but these approaches can require language-specific heuristics \citep{gpt2}, limit out-of-domain performance \citep{sharami2023systematic}, and can affect prompting and truncated sampling in unpredictable ways.\footnote{For example, whether or not a prompt should end in whitespace depends on details of the underlying subwod algorithm used.} \citet{subword} downsamples characters using subword information and has shown promising results in machine translation tasks. The amount of high-frequency information in images and audio means that tokenization cannot be performed losslessly, and instead clustering \citep{hubert} or discrete auto-encoders \citep{dalle} are used to compress the inputs, which lose information and likely limit generative model performance. Our patches are analogous to traditional lossless tokens, and the Local model performs the role of mapping a hidden state to a distribution over possible patches.

\section{Conclusion}
We introduced \textsc{MegaByte}, a scaleable architecture for modeling long sequences.
\textsc{MegaByte} outperforms existing byte-level models across a range of tasks and modalities, allowing large models of sequences of over 1 million tokens. It also gives competitive language modeling results with subword models, which may allow byte-level models to replace tokenization. However, the scale of experiments here is far below those of state-of-the-art language models \citep{gpt3}, and future work should explore scaling \textsc{MegaByte} to much larger models and datasets.

\end{document}
